\documentclass[12pt,a4paper]{report}

% ---------- encoding & fonts ----------
\usepackage[T1]{fontenc}
\usepackage[utf8]{inputenc}
\usepackage{lmodern}

% ---------- geometry ----------
\usepackage[top=1in,bottom=1in,left=1.25in,right=1in]{geometry}

% ---------- spacing ----------
\usepackage{setspace}

% ---------- graphics ----------
\usepackage{graphicx}
\graphicspath{{figures/}}

% ---------- tables ----------
\usepackage{booktabs}
\usepackage{array}
\usepackage{tabularx}
\usepackage{longtable}
\usepackage{multirow}

% ---------- math ----------
\usepackage{amsmath}
\usepackage{amssymb}

% ---------- colors & listings ----------
\usepackage{xcolor}
\usepackage{listings}

% ---------- hyperref (load last among setup pkgs) ----------
\usepackage[hidelinks]{hyperref}
\hypersetup{
  pdfauthor={Ayushi Ranjan, Swadha Kumari, Vraj Soni, Suryanshu Banerjee},
  pdftitle={LeBlanc: Automated Prompt Enrichment and Iterative Repair for Secure LLM-Generated Code}
}

% ---------- headers ----------
\usepackage{fancyhdr}
\pagestyle{fancy}
\fancyhf{}
\fancyhead[L]{\small LeBlanc: Automated Prompt Enrichment and Iterative Repair}
\fancyhead[R]{\small \thepage}
\renewcommand{\headrulewidth}{0.4pt}

% ---------- chapter / section titles ----------
\usepackage{titlesec}
\titleformat{\chapter}[block]
  {\large\bfseries\centering}{Chapter \thechapter:}{0.5em}{}
\titlespacing*{\chapter}{0pt}{20pt}{10pt}

% ---------- bibliography ----------
\usepackage[style=ieee,backend=biber]{biblatex}
\addbibresource{references.bib}

% ---------- tikz ----------
\usepackage{tikz}
\usepackage{pgfplots}
\pgfplotsset{compat=1.18}
\usetikzlibrary{shapes.geometric,arrows.meta,positioning,fit,backgrounds}

% ---------- caption ----------
\usepackage{caption}
\usepackage{subcaption}
\usepackage{float}

% =========================================================
\begin{document}
\onehalfspacing

% =========================================================
% PAGE 1 -- TITLE PAGE
% =========================================================
\begin{titlepage}
\centering
\vspace*{0.2in}

{\Large\bfseries K.J.\ Somaiya School of Engineering}\\[0.08in]
{\large Department of Information Technology}\\[0.04in]
{\normalsize Affiliated to University of Mumbai}\\[0.35in]

\rule{\linewidth}{1.2pt}\\[0.18in]
{\LARGE\bfseries LeBlanc: Automated Prompt Enrichment and\\[0.1in]
Iterative Repair for Secure LLM-Generated Code}\\[0.12in]
{\large A Multi-Model Comparative Study}\\[0.18in]
\rule{\linewidth}{1.2pt}\\[0.28in]

{\normalsize\itshape Project Report submitted in partial fulfillment of the requirements\\[0.04in]
for the Degree of Bachelor of Technology in Information Technology\\[0.04in]
Semester VI, Academic Year 2025--26}\\[0.35in]

{\normalsize\bfseries Submitted by:}\\[0.12in]
\begin{tabular}{@{}ll@{}}
  Ayushi Ranjan       & Roll No.\ XXXX01 \\[0.04in]
  Swadha Kumari       & Roll No.\ XXXX02 \\[0.04in]
  Vraj Soni           & Roll No.\ XXXX03 \\[0.04in]
  Suryanshu Banerjee  & Roll No.\ XXXX04 \\
\end{tabular}\\[0.35in]

{\normalsize\bfseries Project Guide:}\\[0.08in]
{\normalsize Prof.\ Deepti Patole}\\[0.04in]
{\normalsize Department of Information Technology}\\[0.04in]
{\normalsize K.J.\ Somaiya School of Engineering}\\[0.35in]

% Replace with actual logo file if available; comment out if not present
% \includegraphics[width=1.2in]{kjsse_logo.png}\\[0.25in]

{\normalsize K.J.\ Somaiya School of Engineering, Vidyavihar, Mumbai -- 400077}\\[0.06in]
{\normalsize May 2026}
\end{titlepage}

% =========================================================
% PAGE 2 -- CERTIFICATE
% =========================================================
\newpage
\thispagestyle{empty}
\begin{center}
  {\Large\bfseries K.J.\ Somaiya School of Engineering}\\[0.06in]
  {\large Department of Information Technology}\\[0.04in]
  {\normalsize Vidyavihar, Mumbai -- 400077}\\[0.3in]
  {\Large\bfseries CERTIFICATE}
\end{center}

\vspace{0.25in}
\noindent This is to certify that the project titled \textbf{``LeBlanc: Automated Prompt
Enrichment and Iterative Repair for Secure LLM-Generated Code''} has been carried
out by the following students of the Department of Information Technology, K.J.\
Somaiya School of Engineering, in partial fulfillment of the requirements for the
Degree of Bachelor of Technology in Information Technology, University of Mumbai,
during the academic year 2025--26.

\vspace{0.2in}
\begin{center}
\begin{tabular}{ll}
\toprule
\textbf{Name} & \textbf{Roll Number} \\
\midrule
Ayushi Ranjan       & XXXX01 \\
Swadha Kumari       & XXXX02 \\
Vraj Soni           & XXXX03 \\
Suryanshu Banerjee  & XXXX04 \\
\bottomrule
\end{tabular}
\end{center}

\vspace{0.3in}
\noindent The work has been carried out under our supervision and is found to be
satisfactory and worthy of acceptance.

\vspace{0.7in}
\noindent
\begin{tabular}{@{}p{2.8in}p{2.5in}@{}}
  \textbf{Prof.\ Deepti Patole}       & \textbf{Dr.\ [HOD Name]}             \\[0.04in]
  Project Guide                       & Head of Department                   \\[0.04in]
  Dept.\ of Information Technology    & Dept.\ of Information Technology     \\[0.04in]
  K.J.\ Somaiya School of Engineering & K.J.\ Somaiya School of Engineering  \\
\end{tabular}

\vspace{0.5in}
\noindent\textbf{Date:}\enspace\underline{\hspace{1.8in}}\hfill\textbf{Place:}\enspace Mumbai

% =========================================================
% PAGE 3 -- DECLARATION
% =========================================================
\newpage
\thispagestyle{empty}
\begin{center}
  {\Large\bfseries DECLARATION}
\end{center}

\vspace{0.25in}
\noindent We, the undersigned students of the Department of Information Technology,
K.J.\ Somaiya School of Engineering, hereby declare that the project titled
\textbf{``LeBlanc: Automated Prompt Enrichment and Iterative Repair for Secure
LLM-Generated Code''} submitted in partial fulfillment of the requirements for the
Degree of Bachelor of Technology in Information Technology, University of Mumbai,
is our own original work carried out during the academic year 2025--26.

\vspace{0.15in}
\noindent We further declare that:
\begin{enumerate}
  \item This work has not been submitted, in whole or in part, for any other degree
        or examination at this or any other institution.
  \item All sources of information and literature used in this project have been duly
        acknowledged and cited in the references section.
  \item The project does not involve any fabrication, falsification, or
        misrepresentation of data or results.
  \item All experimental pipeline runs, benchmark evaluations, and software
        artifacts described in this report were developed and executed by our team
        during the stated period.
  \item Where third-party datasets, tools, or APIs have been used, their origins
        have been clearly identified and all applicable terms of use have been
        followed.
\end{enumerate}

\vspace{0.45in}
\noindent
\begin{tabular}{@{}p{2.8in}p{2.5in}@{}}
  Ayushi Ranjan      & Swadha Kumari      \\[0.5in]
  Vraj Soni          & Suryanshu Banerjee \\
\end{tabular}

\vspace{0.3in}
\noindent\textbf{Date:}\enspace\underline{\hspace{1.8in}}\hfill\textbf{Place:}\enspace Mumbai

% =========================================================
% PAGE 4 -- ACKNOWLEDGEMENTS
% =========================================================
\newpage
\thispagestyle{empty}
\begin{center}
  {\Large\bfseries ACKNOWLEDGEMENTS}
\end{center}

\vspace{0.25in}
\noindent We thank Prof.\ Deepti Patole for her supervision throughout this project.
Her feedback on metrics design, system architecture, and research framing shaped
LeBlanc into a more rigorous system than we could have built without that input.

\vspace{0.15in}
\noindent We acknowledge the authors of the SecurityEval, LLMSecEval, and SecuCoGen
datasets for making their benchmark prompts and annotations publicly available.
Our empirical evaluation relies on these resources directly.

\vspace{0.15in}
\noindent We thank Semgrep Inc.\ and the Python Security Project for maintaining
the open-source static analysis tools that form the detection core of Engine B.
Without reliable, deterministic scanners the entire pipeline would rest on LLM
opinion rather than rule-based analysis.

\vspace{0.15in}
\noindent API access for Google Gemini 2.5 Flash was provisioned through Google
AI Studio. API access for Cohere Command-A was provisioned through the Cohere
platform, and access for Meta Llama 3.1-8B was provisioned through Groq's
inference API.

\vspace{0.15in}
\noindent Finally, we thank the Department of Information Technology, K.J.\
Somaiya School of Engineering, for providing the infrastructure and academic
environment that made this project possible.

% =========================================================
% PAGE 5 -- ABSTRACT
% =========================================================
\newpage
\thispagestyle{empty}
\begin{center}
  {\Large\bfseries ABSTRACT}
\end{center}

\vspace{0.2in}
\noindent Large language models are increasingly used to generate production code,
yet empirical studies show that a significant fraction of their output contains
exploitable security weaknesses. Existing evaluations are largely single-turn and
single-model: a prompt goes in, code comes out, a scanner marks it vulnerable or
clean, and the loop ends there. No prior system closes the full cycle from
proactively warning the model, to detecting what it still gets wrong, to
autonomously repairing the output---across multiple models running simultaneously.

\vspace{0.15in}
\noindent This report presents \textbf{LeBlanc}, a three-engine pipeline designed
to fill that gap. Engine~A performs rule-based CWE-aware prompt enrichment: before
any LLM call, it parses the user's prompt for security-relevant keywords and
injects structured warnings drawn from a hand-curated CWE mapping file.
Engine~B runs the generated code through Semgrep and Bandit in parallel, mapping
raw scanner output to CWE identifiers, severity levels, and human-readable
vulnerability categories. Engine~C executes an iterative repair loop: when
Engine~B finds vulnerabilities it constructs a structured repair prompt from the
scanner findings and resubmits it to the originating LLM, re-scanning the result
after each attempt for up to three iterations or until the output is clean.

\vspace{0.15in}
\noindent We evaluate the pipeline across three models---Cohere Command-A,
Google Gemini 2.5 Flash, and Meta Llama 3.1-8B---on 33 security-sensitive prompts
drawn from the LLMSecEval benchmark, covering ten CWE types across five
vulnerability categories. Each prompt is run in three modes: plain, CWE-enriched,
and enriched with iterative repair. Results are benchmarked against the GitHub
Copilot 2022 baseline vulnerability rate of 63.4\% reported by Pearce et
al.~\cite{pearce2022asleep}.

\vspace{0.15in}
\noindent Key findings: Command-A and Gemini 2.5 Flash achieve vulnerability rates
of 57.1\% and 54.5\% respectively in plain mode, both below the Copilot baseline.
CWE-aware prompt enrichment reduces Command-A's vulnerability rate by a further
42.9\%. The iterative repair loop resolves vulnerabilities in 77\% of runs that
enter Engine~C, with a mean of 1.4 iterations to convergence. Llama 3.1-8B shows
degraded performance under enrichment; we attribute this to the model's limited
instruction-following capacity at 8B parameters rather than a failure of the
enrichment strategy itself. This finding motivates future evaluation on larger
open-weight models such as Llama 3.3-70B.

\vspace{0.15in}
\noindent The system is implemented as a web application with a Flask backend,
an HTML/JavaScript frontend, and a SQLite database that logs every prompt,
generated code, scan result, and repair iteration for reproducibility. A VS~Code
extension and a full three-repetition dataset collection run are identified as
near-term future work.

\vspace{0.18in}
\noindent\textbf{Keywords:} LLM code generation, software security, static analysis,
prompt engineering, CWE, iterative repair, Semgrep, Bandit, vulnerability rate,
multi-model evaluation

% =========================================================
% FRONT MATTER -- TOC / LOF / LOT
% =========================================================
\newpage
\tableofcontents

\newpage
\listoffigures

\newpage
\listoftables

\end{document}
